% GENERATED FILE. Edit canonical lessons, then run npm run generate:books.
% canonical-source-hash: fnv1a64:982adcae68a2a642
% canonical-lessons: GE-C491-halbbruder, GE-C491-pate, GE-C491-patenkind, GE-C491-ehepartner, GE-C491-nachwuchs

\chapter{Paten und Nachwuchs}
\label{ch:ge-paten-und-nachwuchs}
\hlchaptermodality{491}

\begin{chapteropening}
\textbf{By the end of this chapter:} \emph{I can talk about godparents, partners and babies.}

Complete the last lesson of chapter 491: I can talk about godparents, partners and babies.
\end{chapteropening}

\section[der Halbbruder]{der Halbbruder — half-brother}

\label{lesson:GE-C491-halbbruder}



\begin{quote}
Before the new one: say the German for twin, then the German for only child.
\end{quote}

\subsection*{You'll want to know: der Halbbruder}
\textbf{der Halbbruder} — \textquotedblleft{}half-brother\textquotedblright{}.

\begin{quote}
Mein Halbbruder ist zehn Jahre jünger als ich.
\end{quote}

\begin{quote}
\emph{My half-brother is ten years younger than me.}
\end{quote}

Plural: \emph{die Halbbrüder}. Likewise \emph{die Halbschwester}.

\begin{cousinweb}
\emph{Halbbruder} is built on \emph{halb} \textquotedblleft{}half\textquotedblright{}, inherited from Old High German \emph{halb}, from Proto-Germanic \emph{halbaz} (EtymDB), and \emph{Bruder} \textquotedblleft{}brother\textquotedblright{}.
\end{cousinweb}

\begin{grammarlens}[title={What you've built}]
One more word for telling, explaining and describing at B1.
\end{grammarlens}

\subsection*{Your turn}
\begin{itemize}
\raggedright
  \item \emph{Say it:} \emph{der Halbbruder}
  \item \emph{Say it:} \emph{der Halbbruder}, once more
  \item \emph{Say it:} a sentence of your own with \emph{der Halbbruder}
  \item \emph{Recall:} say the German for twin, then the German for only child, then say \emph{der Halbbruder} again
\end{itemize}

\begin{tcolorbox}[breakable,colback=teal!4,colframe=teal!35!black,title={Before you move on}]
What does der Halbbruder mean? (\textquotedblleft{}Half-brother\textquotedblright{}.) Say it once more.
\end{tcolorbox}
\section[der Pate]{der Pate — godfather; godparent}

\label{lesson:GE-C491-pate}



\begin{quote}
Before the new one: say the German for only child, then the German for half-brother.
\end{quote}

\subsection*{You'll want to know: der Pate}
\textbf{der Pate} — \textquotedblleft{}godfather; godparent\textquotedblright{}.

\begin{quote}
Mein Onkel ist der Pate meiner Tochter.
\end{quote}

\begin{quote}
\emph{My uncle is my daughter's godfather.}
\end{quote}

Weak noun: \emph{den Paten}. A godmother is \emph{die Patin} or \emph{die Patentante}.

\begin{cousinweb}
\emph{Pate} comes from Middle High German \emph{pate} (EtymDB), usually traced to Church Latin \emph{pater} \textquotedblleft{}father\textquotedblright{}, as in \emph{pater spiritualis} \textquotedblleft{}spiritual father\textquotedblright{}.
\end{cousinweb}

\begin{grammarlens}[title={What you've built}]
One more word for telling, explaining and describing at B1.
\end{grammarlens}

\subsection*{Your turn}
\begin{itemize}
\raggedright
  \item \emph{Say it:} \emph{der Pate}
  \item \emph{Say it:} \emph{der Pate}, once more
  \item \emph{Say it:} a sentence of your own with \emph{der Pate}
  \item \emph{Recall:} say the German for only child, then the German for half-brother, then say \emph{der Pate} again
\end{itemize}

\begin{tcolorbox}[breakable,colback=teal!4,colframe=teal!35!black,title={Before you move on}]
What does der Pate mean? (\textquotedblleft{}Godfather; godparent\textquotedblright{}.) Say it once more.
\end{tcolorbox}
\section[das Patenkind]{das Patenkind — godchild}

\label{lesson:GE-C491-patenkind}



\begin{quote}
Before the new one: say the German for half-brother, then the German for godfather.
\end{quote}

\subsection*{You'll want to know: das Patenkind}
\textbf{das Patenkind} — \textquotedblleft{}godchild\textquotedblright{}.

\begin{quote}
Zum Geburtstag schenke ich meinem Patenkind ein Buch.
\end{quote}

\begin{quote}
\emph{I'm giving my godchild a book for their birthday.}
\end{quote}

Plural: \emph{die Patenkinder}.

\begin{cousinweb}
\emph{Patenkind} is built on \emph{Pate} \textquotedblleft{}godparent\textquotedblright{}, from Middle High German \emph{pate} (EtymDB), and \emph{Kind} \textquotedblleft{}child\textquotedblright{}.
\end{cousinweb}

\begin{grammarlens}[title={What you've built}]
One more word for telling, explaining and describing at B1.
\end{grammarlens}

\subsection*{Your turn}
\begin{itemize}
\raggedright
  \item \emph{Say it:} \emph{das Patenkind}
  \item \emph{Say it:} \emph{das Patenkind}, once more
  \item \emph{Say it:} a sentence of your own with \emph{das Patenkind}
  \item \emph{Recall:} say the German for half-brother, then the German for godfather, then say \emph{das Patenkind} again
\end{itemize}

\begin{tcolorbox}[breakable,colback=teal!4,colframe=teal!35!black,title={Before you move on}]
What does das Patenkind mean? (\textquotedblleft{}Godchild\textquotedblright{}.) Say it once more.
\end{tcolorbox}
\section[der Ehepartner]{der Ehepartner — spouse}

\label{lesson:GE-C491-ehepartner}



\begin{quote}
Before the new one: say the German for godfather, then the German for godchild.
\end{quote}

\subsection*{You'll want to know: der Ehepartner}
\textbf{der Ehepartner} — \textquotedblleft{}spouse\textquotedblright{}.

\begin{quote}
Auf dem Formular steht auch der Name des Ehepartners.
\end{quote}

\begin{quote}
\emph{The spouse's name is also on the form.}
\end{quote}

Plural: \emph{die Ehepartner}. A woman is \emph{die Ehepartnerin}.

\begin{cousinweb}
\emph{Ehepartner} is built on \emph{Ehe} \textquotedblleft{}marriage\textquotedblright{} and \emph{Partner}, which was borrowed from English \emph{partner} (EtymDB).
\end{cousinweb}

\begin{grammarlens}[title={What you've built}]
One more word for telling, explaining and describing at B1.
\end{grammarlens}

\subsection*{Your turn}
\begin{itemize}
\raggedright
  \item \emph{Say it:} \emph{der Ehepartner}
  \item \emph{Say it:} \emph{der Ehepartner}, once more
  \item \emph{Say it:} a sentence of your own with \emph{der Ehepartner}
  \item \emph{Recall:} say the German for godfather, then the German for godchild, then say \emph{der Ehepartner} again
\end{itemize}

\begin{tcolorbox}[breakable,colback=teal!4,colframe=teal!35!black,title={Before you move on}]
What does der Ehepartner mean? (\textquotedblleft{}Spouse\textquotedblright{}.) Say it once more.
\end{tcolorbox}
\section[der Nachwuchs]{der Nachwuchs — offspring, new arrival; young talent}

\label{lesson:GE-C491-nachwuchs}



\begin{quote}
Before the new one: say the German for godchild, then the German for spouse.
\end{quote}

\subsection*{You'll want to know: der Nachwuchs}
\textbf{der Nachwuchs} — \textquotedblleft{}offspring, new arrival; young talent\textquotedblright{}.

\begin{quote}
Unsere Nachbarn haben Nachwuchs bekommen.
\end{quote}

\begin{quote}
\emph{Our neighbours have had a baby.}
\end{quote}

No plural. Also young talent: \emph{der Nachwuchs im Fußball}.

\begin{cousinweb}
\emph{Nachwuchs} is built on \emph{nach} \textquotedblleft{}after\textquotedblright{} and \emph{Wuchs} \textquotedblleft{}growth\textquotedblright{}, from \emph{wachsen} \textquotedblleft{}to grow\textquotedblright{}, inherited from Old High German \emph{wahsan}, from Proto-Germanic \emph{wahsijaną} (EtymDB).
\end{cousinweb}

\begin{grammarlens}[title={What you've built}]
One more word for telling, explaining and describing at B1.
\end{grammarlens}

\subsection*{Your turn}
\begin{itemize}
\raggedright
  \item \emph{Say it:} \emph{der Nachwuchs}
  \item \emph{Say it:} \emph{der Nachwuchs}, once more
  \item \emph{Say it:} a sentence of your own with \emph{der Nachwuchs}
  \item \emph{Recall:} say the German for godchild, then the German for spouse, then say \emph{der Nachwuchs} again
\end{itemize}

\begin{tcolorbox}[breakable,colback=teal!4,colframe=teal!35!black,title={Before you move on}]
What does der Nachwuchs mean? (\textquotedblleft{}Offspring, new arrival; young talent\textquotedblright{}.) Say it once more.
\end{tcolorbox}
